\documentclass[11pt,a4paper]{article}
\usepackage[T1]{fontenc}
\usepackage[utf8]{inputenc}
\usepackage{lmodern}
\usepackage[margin=27mm,headheight=14pt]{geometry}
\usepackage{amsmath,amssymb,amsthm,mathtools}
\usepackage{microtype,booktabs,array,longtable}
\usepackage{enumitem,xurl,fancyhdr}
\usepackage[hidelinks]{hyperref}
\usepackage[nameinlink,noabbrev]{cleveref}
\hypersetup{pdftitle={The Lacunarity Boundary: Two-Moment Rigidity and Sharp Recovery of Geometric Uniform Spectra},pdfauthor={Research manuscript prepared with ChatGPT},pdfsubject={Critical dyadic recovery, moment rigidity, sharp Holder stability, and testing complexity}}
\pagestyle{fancy}\fancyhf{}
\fancyhead[L]{\small The lacunarity boundary}
\fancyhead[R]{\small Rigidity, recovery, and testing}
\fancyfoot[C]{\thepage}
\setlength{\parskip}{3pt}
\setlength{\emergencystretch}{3em}
\setlist[enumerate]{itemsep=4pt,topsep=4pt}
\numberwithin{equation}{section}
\newtheorem{theorem}{Theorem}[section]
\newtheorem{lemma}[theorem]{Lemma}
\newtheorem{proposition}[theorem]{Proposition}
\newtheorem{corollary}[theorem]{Corollary}
\theoremstyle{definition}
\newtheorem{definition}[theorem]{Definition}
\newtheorem{question}{Research question}
\theoremstyle{remark}
\newtheorem{remark}[theorem]{Remark}
% ed. (2026-09-30): unnumbered environment for ProveIt editorial notes (the theorem counter is unchanged).
\newtheorem*{ednote}{Editorial note (ProveIt, 2026-09-30)}
\DeclareMathOperator{\TV}{TV}
\DeclareMathOperator{\Law}{Law}
\DeclareMathOperator{\Var}{Var}
\DeclareMathOperator{\supp}{supp}
\DeclareMathOperator{\sinc}{sinc}
\newcommand{\R}{\mathbb R}
\newcommand{\N}{\mathbb N}
\newcommand{\E}{\mathbb E}
\newcommand{\PP}{\mathbb P}
\newcommand{\C}{\mathcal C}
\newcommand{\K}{\mathcal K}
\newcommand{\eps}{\varepsilon}
\newcommand{\dd}{\,\mathrm d}
\newcommand{\norm}[1]{\left\lVert #1\right\rVert}
\newcommand{\dK}{d_{\mathrm K}}
\newcommand{\repoSHA}{f4457a1d19701d32e8f324240c497653d235b53e}
\newcommand{\blobSHA}{50cc46f4084386b48378cbc9c1d942c66ce64741}

\title{\textbf{The Lacunarity Boundary}\\[7pt]
\Large Two-Moment Rigidity and Sharp Recovery\\of Geometric Uniform Spectra\\[9pt]
\large Resolving the critical dyadic separation question,\\with sharp prefix conditioning and quartic testing complexity}
\author{Research manuscript prepared for Vladimir Reshetnikov\\
\small AI-assisted mathematical research with conventional proofs}
\date{30 September 2026}

\begin{document}
\maketitle
\begin{abstract}
Let $X_a=\sum_{j\ge1}a_jU_j$, where the $U_j$ are independent uniforms
on $[-1,1]$. We resolve the exact separation-boundary question posed in the
ProveIt article \emph{Anchored Recovery of the Dyadic Uniform Spectrum}:
under $a_{j+1}\le a_j/2$ and $\sum a_j^2=1/3$, the full coefficient
sequence is recoverable from its law with the sharp exponent $1/2$.
In fact, the fourth moment alone gives the global certificate
\[
 \sup_j|a_j-2^{-j}|^2
 \le \frac{75}{4}\left(\frac{19}{675}-\E X_a^4\right).
\]
The expression in parentheses is nonnegative and vanishes only at the exact
dyadic spectrum. Thus two moments identify an infinite spectrum inside this
one-sided constraint class. A deleted-factor construction proves matching
square-root lower bounds in total variation and Kolmogorov distance, even
inside the predecessor's precise infinitely smooth class. Fixed-prefix local
Lipschitz constants grow as $\Theta(2^r)$, and testing the reference against
spectral distance at least $\delta$ requires and suffices to use
$\Theta(\delta^{-4})$ observations. The statistical lower bound uses a
boundary-sensitive Hellinger estimate, not a product-total-variation shortcut.
We also prove exact support rigidity, extend the method to every geometric
reference and to nongeometric relative-ratio cones, and separate these results
from the still-open pairwise and small-slack problems.
\end{abstract}

\noindent\textbf{Status and scope.} The new statements are proved in this
manuscript, but have not been checked in Lean or independently peer reviewed.
The claim of novelty is a proposed repository-relative research contribution,
not an exhaustive literature-priority determination. The basic summation-by-parts,
convolution, and concentration tools are classical.

\clearpage
\begingroup\small\setlength{\parskip}{0pt}\tableofcontents\endgroup
\clearpage

\section{The repository question and the main conclusions}
\label{sec:introduction}

\subsection{Why the exact boundary is a different inverse problem}

The ProveIt Fabius-function collection contains the research article
\emph{Anchored Recovery of the Dyadic Uniform Spectrum} \cite{anchored}.
Its third further-research question asks what happens under the exact inequality
$a_{j+1}\le a_j/2$, with variance fixed at $1/9$. That article constructs
very small law perturbations when any fixed slack above $1/2$ is available,
but explicitly leaves the boundary case unresolved. Its construction preserves
many power sums by compensating changes of neighboring scales; the exact
one-sided constraint obstructs those changes.

The present answer is positive and considerably more explicit than a qualitative
inverse-continuity statement. At the boundary, the fourth moment is an exposing
functional for the reference spectrum: its deficit controls every squared
coefficient. This reduces an infinite-dimensional anchored inverse problem to
one scalar measurement. It does \emph{not} say that two moments determine all
spectra, or that two freely varying spectra admit the same inverse estimate.
The reference is the unique moment extremizer inside the specified cone.

For comparison, the predecessor proves the anchored scale
$\exp[-\Theta(\sqrt{\log(1/\eps)})]$ on its smooth classes with every fixed
ratio bound $\rho>1/2$ \cite{anchored}. Our critical-boundary proofs do not
import that unformalized quantitative result. The comparison with it is an
attributed repository result; the positive and sharp boundary theorems below
are established independently.

\subsection{The dyadic statement}

Throughout, $U_1,U_2,\ldots$ are independent random variables uniform on
$[-1,1]$. Put
\begin{equation}
 a_j^0=2^{-j},\qquad X_0=\sum_{j\ge1}2^{-j}U_j,
 \qquad \mu_0=\Law(X_0).
\end{equation}
The density of $\mu_0$ is the centered Rvachev up-function, here normalized
on $[-1,1]$. This probability interpretation and the compactly supported smooth
function are established background \cite{arias}; we derive the moment and
norm formulas needed below directly.

\begin{definition}[Critical dyadic class]
\label{def:critical}
Let
\begin{equation}
 \C_* =\left\{a=(a_j)_{j\ge1}:a_j\ge0,\quad
 a_{j+1}\le\tfrac12a_j,\quad \sum_{j\ge1}a_j^2=\tfrac13\right\}.
 \label{eq:critical-class}
\end{equation}
For $a\in\C_*$ write $X_a=\sum a_jU_j$, $\mu_a=\Law(X_a)$,
$m_k(a)=\E X_a^k$, and
\[
 d_\infty(a,b)=\sup_j|a_j-b_j|,\qquad
 d_r(a,b)=\max_{1\le j\le r}|a_j-b_j|.
\]
Finite sequences are padded by zeros. All such random series converge
absolutely because $a_j\le a_1 2^{1-j}$.
\end{definition}

Our total variation convention is
$\TV(\mu,\nu)=\sup_B|\mu(B)-\nu(B)|$, one-half the $L^1$ distance for
densities. Kolmogorov distance is
$\dK(\mu,\nu)=\sup_x|F_\mu(x)-F_\nu(x)|$.
In particular $\dK\le\TV$.

\begin{theorem}[Critical dyadic recovery]
\label{thm:dyadic-intro}
Every $a\in\C_*$ is supported probabilistically in $[-1,1]$, has
$m_2(a)=1/9$, and satisfies
\begin{align}
 e_4(a)&:=\frac{19}{675}-m_4(a)\ge0, \label{eq:e4-dyadic}\\
 d_\infty(a,a^0)^2&\le\frac{75}{4}e_4(a), \label{eq:dyadic-moment-main}\\
 d_r(a,a^0)&\le\frac{75}{4}2^r e_4(a)\qquad(r\ge1).
 \label{eq:dyadic-prefix-main}
\end{align}
Equality $e_4(a)=0$ holds if and only if $a=a^0$. Consequently,
\begin{align}
 d_\infty(a,a^0)&\le\sqrt{\tfrac{75}{4}\TV(\mu_a,\mu_0)},
 \label{eq:dyadic-TV-main}\\
 d_\infty(a,a^0)&\le\sqrt{\tfrac{75}{2}\dK(\mu_a,\mu_0)}.
 \label{eq:dyadic-K-main}
\end{align}
The exponent $1/2$ is optimal for both distances, even after imposing the
exact smooth-class conditions in \cref{def:smooth}.
\end{theorem}

The two distinct scales in \eqref{eq:dyadic-moment-main} and
\eqref{eq:dyadic-prefix-main} are important. An error concentrated at a
coefficient of size $h$ changes its square by order $h^2$, so a global loss
naturally has a square root. A fixed coefficient is bounded away from zero,
so dividing a squared-coefficient error by that coefficient gives a Lipschitz
bound. We prove that its dependence on the prefix length is also sharp.

\begin{theorem}[Statistical and conditioning consequences]
\label{thm:stat-intro}
At the dyadic reference, the local Lipschitz constant for the first $r$
coefficients is $\Theta(2^r)$ in either law distance. Testing
$a=a^0$ against $a\in\C_*$ with $d_\infty(a,a^0)\ge\delta$, with both
error probabilities at most $1/4$, has sample complexity
$\Theta(\delta^{-4})$. Both lower bounds remain valid in the exact smooth
subclass. An honest moment-based confidence set over $\C_*$ contracts at
$a^0$ with radius $O(N^{-1/4})$, and its fixed-prefix radius is
$O(2^rN^{-1/2})$.
\end{theorem}

These are anchored testing and confidence statements. In particular, no
claim is made about a globally optimal estimator between arbitrary spectra.
At a single known point the constant estimator has zero loss, so a bare
``pointwise minimax rate'' would not express the content of the testing theorem.

\subsection{Established background and the proposed contribution}

Billey and Swanson study generalized uniform sums and their unique
sequence-space parametrization, including a topology allowing Gaussian
components \cite{billey}. Their qualitative identifiability is not a new
claim here. Arias de Reyna develops the smooth compactly supported dyadic
function and its probability and derivative representations \cite{arias}.
Hoeffding's bounded-variable inequality supplies a standard concentration
input \cite{hoeffding}; a short proof of the version used here is included.

The proposed contribution is the quantitative rigidity at the exact
one-sided ratio boundary, together with smooth sharpness witnesses, sharp
prefix conditioning, Hellinger-controlled testing complexity, and the
relative-ratio extension. The algebraic core is elementary enough to isolate
for formalization. Elementary tools do not make the inverse conclusion
automatic: the square-slack bound, preservation of the predecessor's
infinitely many smoothness inequalities, and the fourth-order Hellinger
estimate each address a separate obstruction.

\section{A geometric majorization identity}
\label{sec:majorization}

Fix $0<\rho<1$, write $q=\rho^2$, and fix $S>0$. Define
\begin{equation}
 \C_{\rho,S}=\left\{a_j\ge0:a_{j+1}\le\rho a_j,
                     \ \sum a_j^2=S\right\},
 \qquad
 g_j=\sqrt{S(1-q)}\,\rho^{j-1}.
 \label{eq:general-class}
\end{equation}
Thus $g$ is the geometric reference with exactly the same ratio as the
constraint. When $\rho>1/2$, this reference is \emph{not} the dyadic reference;
our general theorem therefore does not conflict with the predecessor's
instability at the dyadic anchor under a looser ratio bound.

Set
\begin{equation}
 p_j=\frac{a_j^2}{S},\quad b_j=(1-q)q^{j-1},\quad
 c_q=\frac{1-q}{1+q},\quad
 \Delta=\sum_jp_j^2-c_q,
 \quad d=p_1-(1-q).
 \label{eq:normalized}
\end{equation}
Both $p$ and $b$ are probability sequences and $p_{j+1}\le qp_j$.

\begin{lemma}[Majorization with an exact remainder]
\label{lem:energy}
For every probability sequence $p$ satisfying $p_{j+1}\le qp_j$, put
$D_k=\sum_{j\le k}(p_j-b_j)$. Then $D_k\ge0$ and
\begin{equation}
 \boxed{\quad
 \sum_jp_j^2-c_q
 =\sum_j(p_j-b_j)^2
   +2(1-q)^2\sum_{k\ge1}q^{k-1}D_k.\quad}
 \label{eq:energy}
\end{equation}
In particular $\sum p_j^2\ge c_q$, with equality only for $p=b$.
\end{lemma}
\begin{proof}
Iterating the ratio inequality gives $p_{j+k}\le q^kp_j$. Summing over
$j\ge1$ yields $\sum_{j>k}p_j\le q^k$. Therefore
\[
 D_k=q^k-\sum_{j>k}p_j\ge0.
\]
Since $\sum b_j^2=c_q$, expanding the square gives
\[
 \sum p_j^2-c_q=\sum(p_j-b_j)^2+2\sum b_j(p_j-b_j).
\]
Finite summation by parts, followed by a limit, transforms the last sum into
$\sum_{k\ge1}(b_k-b_{k+1})D_k$. The boundary term is
$b_ND_N\to0$, because $D_N$ is bounded and $b_N\to0$.
All series converge absolutely; in particular $0\le D_k\le1$ and the weights
are geometric. Finally $b_k-b_{k+1}=(1-q)^2q^{k-1}$, proving the identity.
If its left side vanishes, the squared-error term is zero coordinatewise.
\end{proof}

The next refinement is stronger than discarding the nonnegative remainder in
\eqref{eq:energy}. A direct $\ell^2$ estimate followed by taking square roots
would lose an unnecessary additional square root. We instead first control
the head excess and then all envelope slacks.

\begin{lemma}[Sharp head-excess inequality]
\label{lem:head}
With the notation \eqref{eq:normalized}, $0\le d\le q$ and
\begin{equation}
 \Delta\ge 2c_qd+\frac{2}{1+q}d^2
          =\frac{2}{1+q}d(1-q+d).
 \label{eq:head-quadratic}
\end{equation}
This bound is sharp for every $d\in[0,q]$. Consequently
\begin{equation}
 d\le\mathcal D_q(\Delta):=
 \frac{\sqrt{(1-q)^2+2(1+q)\Delta}-(1-q)}{2}
 \le\frac{\Delta}{2c_q}.
 \label{eq:head-inverse}
\end{equation}
\end{lemma}
\begin{proof}
The bound $p_1\ge1-q$ follows either from $D_1\ge0$ or from
$1\le p_1/(1-q)$. The remaining mass is $1-p_1$. Apply
\cref{lem:energy} to the tail normalized by this mass to obtain
\[
 \sum_{j\ge2}p_j^2\ge c_q(1-p_1)^2.
\]
When the mass is zero the same inequality is immediate. Substitution of
$p_1=1-q+d$ and expansion gives \eqref{eq:head-quadratic}.
Equality is attained by the sequence
\[
 p_1=1-q+d,\qquad
 p_j=(q-d)(1-q)q^{j-2}\quad(j\ge2).
\]
Its first ratio satisfies $p_2\le qp_1$ precisely because $p_1\ge1-q$;
the later ratios equal $q$. Solving the quadratic proves the first bound in
\eqref{eq:head-inverse}, and dropping its positive quadratic term proves
the second.
\end{proof}

\begin{lemma}[Every squared coefficient is controlled by the head]
\label{lem:slack}
For $a\in\C_{\rho,S}$,
\begin{equation}
 |a_j^2-g_j^2|\le Sd\qquad(j\ge1).
 \label{eq:square-slack}
\end{equation}
\end{lemma}
\begin{proof}
Define the nonnegative envelope slack
$R_j=Sp_1q^{j-1}-a_j^2$. The ratio constraint implies
$R_{j+1}\ge qR_j$, and summation gives
\[
 \sum_jR_j=\frac{Sp_1}{1-q}-S=\frac{Sd}{1-q}.
\]
On the other hand, $\sum_{k\ge0}R_{j+k}\ge R_j/(1-q)$, so
$0\le R_j\le Sd$. Since
\[
 a_j^2-g_j^2=Sdq^{j-1}-R_j,
\]
both terms on the right lie in $[0,Sd]$. Their difference has absolute value
at most $Sd$.
\end{proof}

\section{Moment certificates, support rigidity, and global recovery}
\label{sec:recovery}

\subsection{Moments of a uniform sum}

\begin{lemma}
\label{lem:moments}
For $a\in\C_{\rho,S}$, put $T=\sum a_j^4$. Then
\begin{equation}
 m_2(a)=\frac S3,\qquad
 m_4(a)=\frac{S^2}{3}-\frac{2T}{15}.
 \label{eq:moments}
\end{equation}
For the geometric reference $T(g)=S^2c_q$, and hence
\begin{equation}
 e_4(a):=m_4(g)-m_4(a)=\frac{2S^2}{15}\Delta\ge0.
 \label{eq:moment-deficit}
\end{equation}
\end{lemma}
\begin{proof}
For a centered uniform $U$, $\E U^2=1/3$ and $\E U^4=1/5$.
Independence and symmetry give, for a finite sum,
\[
 \E X^4=\frac15\sum a_j^4+\frac23\sum_{i<j}a_i^2a_j^2
        =\frac13\Bigl(\sum a_j^2\Bigr)^2-\frac2{15}\sum a_j^4.
\]
The infinite series is bounded by $\sum a_j\le\sqrt S/(1-\rho)$,
so dominated convergence passes these formulas to the limit.
The geometric series for $T(g)$ and \cref{lem:energy} finish the proof.
\end{proof}

\begin{theorem}[Quantitative geometric rigidity]
\label{thm:geometric-rigidity}
For every $a\in\C_{\rho,S}$,
\begin{align}
 d_\infty(a,g)^2
 &\le S\mathcal D_q(\Delta)
 \le\frac{S}{2c_q}\Delta
 =\frac{15}{4c_qS}e_4(a), \label{eq:general-rigidity}\\
 |a_j-g_j|
 &\le\frac{S\mathcal D_q(\Delta)}{g_j}
 \le\frac{15}{4c_qSg_j}e_4(a). \label{eq:general-prefix}
\end{align}
The intersection of $\C_{\rho,S}$ with the fourth-moment fibre
$m_4(a)=m_4(g)$ is the singleton $\{g\}$.
\end{theorem}
\begin{proof}
For nonnegative $x,y$, $(\sqrt x-\sqrt y)^2\le|x-y|$.
Apply this to \eqref{eq:square-slack} and use \cref{lem:head}.
For the coordinate estimate, use
$|a_j-g_j|=|a_j^2-g_j^2|/(a_j+g_j)\le Sd/g_j$.
If $e_4=0$, \cref{lem:energy} forces $p=b$; nonnegativity then gives $a=g$.
The converse follows by substitution.
\end{proof}

Equivalently, the geometric reference uniquely maximizes kurtosis within the
fixed-variance cone, because
\[
 \frac{m_4(a)}{m_2(a)^2}=3-\frac65\frac{\sum a_j^4}{S^2}.
\]
This is a constrained extremal statement, not a generic finite-moment
identifiability assertion.

\begin{proposition}[The same exponent for summable coefficient losses]
\label{prop:l1}
Write $p_1=1-q+d$. Then
\begin{equation}
 \norm{a-g}_1\le
 \frac{2}{1-\rho}
 \left(\frac{p_1}{1-q}\right)^{1/4}\sqrt{Sd}.
 \label{eq:l1}
\end{equation}
Consequently every $\ell^p$ coefficient loss, $1\le p\le\infty$, is
$O_{\rho,S}(\sqrt{e_4(a)})$.
\end{proposition}
\begin{proof}
Put $A=\sqrt{S(1-q)}$, $B=\sqrt{Sp_1}$, and $Z=Sd$.
The case $Z=0$ follows from rigidity. Otherwise, both the envelope and
\eqref{eq:square-slack} give
\[
 |a_j-g_j|\le\min\{B\rho^{j-1},\ Z/(A\rho^{j-1})\}.
\]
Split the geometric sequence at $\rho^{j-1}\approx t$, where
$t=\sqrt{Z/(AB)}$. If $t<1$, the sum on either side of this crossing is
at most $\sqrt{ZB/A}/(1-\rho)$: on the head use the increasing reciprocal
geometric bound, and on the tail the decreasing direct bound. If $t\ge1$,
the direct bound alone is at most the same quantity. This proves
\eqref{eq:l1}. Since $p_1\le1$, \cref{lem:head} makes its constant uniform.
Finally $\norm{x}_p\le\norm{x}_1^{1/p}\norm{x}_\infty^{1-1/p}$.
\end{proof}

\subsection{An exact support endpoint cannot also move}

\begin{theorem}[Support extremality and rigidity]
\label{thm:support}
For $a\in\C_{\rho,S}$,
\begin{equation}
 \sum_ja_j\le B_{\rho,S}:=
 \sqrt{\frac{S(1+\rho)}{1-\rho}},
 \label{eq:support-bound}
\end{equation}
with equality if and only if $a=g$. The support of $\mu_a$ is exactly
$[-\sum a_j,\sum a_j]$.
\end{theorem}
\begin{proof}
Set $L=\sum a_j>0$. Apply \cref{lem:energy} to $a_j/L$ with ratio
parameter $\rho$, rather than to $a_j^2/S$ with parameter $q$.
It yields
$S/L^2\ge(1-\rho)/(1+\rho)$, with equality exactly for a geometric
sequence of ratio $\rho$. The fixed squared sum then determines its scale.

The support containment is immediate from $|U_j|\le1$. Each finite uniform
sum has its entire Minkowski-sum interval as support. Zero belongs to the
support of any tail: choose a sufficiently long finite part near zero and
bound the remaining deterministic tail length. Thus every point strictly
between $-L$ and $L$ is in the support of the full sum, by using a sufficiently
long head whose support contains that point. Closure includes the endpoints.
\end{proof}

At the dyadic boundary this says $\sum a_j\le1$. Therefore the simultaneous
requirements $a_{j+1}\le a_j/2$, $\sum a_j^2=1/3$, and $\sum a_j=1$
leave \emph{no nontrivial perturbations at all}. This resolves the
exact-support issue at the critical boundary only. It does not resolve the
predecessor's exact-support construction problem when slack $\rho>1/2$
is allowed.

\subsection{From moment error to law error}

If two probability measures are supported in $[-B,B]$, then
\begin{equation}
 |m_4(\mu)-m_4(\nu)|\le B^4\TV(\mu,\nu),\qquad
 |m_4(\mu)-m_4(\nu)|\le2B^4\dK(\mu,\nu).
 \label{eq:moment-metrics}
\end{equation}
For the first inequality, $x^4/B^4$ lies in $[0,1]$ and the layer-cake
representation bounds its expectation difference by total variation. For the
second, Stieltjes integration by parts gives
$-\int_{-B}^B4x^3(F_\mu-F_\nu)\dd x$; the boundary terms vanish when one
first integrates on a slightly larger common interval. The integral of
$4|x|^3$ is $2B^4$.

For $q=1/4$, $S=1/3$, we have $c_q=3/5$, $B_{\rho,S}=1$, and
\begin{equation}
 m_4(a^0)=\frac1{27}-\frac2{225}=\frac{19}{675}.
\end{equation}
% ed. (2026-09-30): Lean counterpart of the reference fourth moment, not cited by the article
\begin{ednote}
The law of $X_{a^0}=\sum_{j\ge1}2^{-j}U_j$ is Rvachev's $\mathrm{up}$ law,
and its fourth moment $19/675$ is machine-checked as
\path{Fabius.upMoment_four}, derived from the rational moment value
\path{Fabius.moment_two}, both in
\path{Analysis/FabiusFunction/Lean/FabiusFunction/OrthogonalPolynomialJacobi.lean}.
No result of this article is formalized.
\end{ednote}
Substituting in \cref{thm:geometric-rigidity} and
\eqref{eq:moment-metrics} proves all upper estimates in
\cref{thm:dyadic-intro}.

For $\rho<1/2$ the dyadic reference is not in $\C_{\rho,1/3}$.
There is in fact a positive separation, not just a failure of membership:
\begin{equation}
 \TV(\mu_a,\mu_0)\ge\eta_\rho,
 \qquad \dK(\mu_a,\mu_0)\ge\tfrac12\eta_\rho,
 \quad
 \eta_\rho=\frac2{135}\left(\frac{1-\rho^2}{1+\rho^2}-\frac35\right)>0.
 \label{eq:subcritical-separation}
\end{equation}
Indeed the fourth-power minimum is larger than the dyadic minimum by the
quantity displayed, and all supports still lie in $[-1,1]$.

\subsection{When the variance is measured rather than fixed}

Suppose $a_{j+1}\le\rho a_j$, but $S=\sum a_j^2$ is not prescribed.
For exact moment measurements with $m_2>0$,
\begin{equation}
 S=3m_2,\qquad T=\frac{15}{2}(3m_2^2-m_4),\qquad
 \Delta=\frac{T}{S^2}-c_q.
 \label{eq:observable-normalization}
\end{equation}
Let $g(S_0)$ be a fixed geometric reference. The triangle inequality and
\cref{thm:geometric-rigidity} give the finite certificate
\begin{equation}
 d_\infty(a,g(S_0))\le
 \sqrt{S\mathcal D_q(\Delta)}
 +\sqrt{1-q}\,|\sqrt S-\sqrt{S_0}|.
 \label{eq:unknown-variance}
\end{equation}
Under a common support bound, law convergence to $\mu_{g(S_0)}$ controls
both moments. Near $S_0>0$, the rational expression in
\eqref{eq:observable-normalization} is locally Lipschitz, so
$|S-S_0|=O(D)$ and $0\le\Delta=O(D)$ for $D=\TV$ or $\dK$.
Equation \eqref{eq:unknown-variance} therefore retains the exponent $1/2$.
With noisy moments, one must propagate measurement intervals; a negative
plug-in value of $\Delta$ is not itself a valid certificate.

\section{Smooth sharpness witnesses inside the original class}
\label{sec:witnesses}

\subsection{Exact derivative norms for separated supports}

\begin{lemma}[Separated-uniform derivative norms]
\label{lem:derivative-norms}
Let all $a_j$ be positive, $\sum a_j<\infty$, and assume
$a_j\ge\sum_{k>j}a_k$ for every $j$. Then $X_a$ has a density
$f_a\in C_c^\infty(\R)$ and
\begin{equation}
 \norm{f_a^{(r)}}_\infty
 =\frac{1}{2^{r+1}\prod_{j=1}^{r+1}a_j}
 \qquad(r=0,1,2,\ldots).
 \label{eq:derivative-norms}
\end{equation}
The hypothesis holds whenever $a_{j+1}\le a_j/2$.
\end{lemma}
\begin{proof}
The characteristic function is $\prod_j\sinc(a_jt)$, with
$\sinc u=\sin(u)/u$. Each fixed positive factor contributes a bound
$\min(1,1/(a_j|t|))$. Using arbitrarily many factors gives arbitrarily
fast polynomial decay, so Fourier inversion gives a $C^\infty$ density.
Its support is compact by absolute summability.

The derivative of the uniform density on $[-a,a]$, in the sense of
distributions, is $(\delta_{-a}-\delta_a)/(2a)$. Differentiating each of
the first $r$ factors once expresses $f_a^{(r)}$ as a signed sum of translates
of the remaining tail density, with coefficient magnitude
$1/\prod_{j\le r}(2a_j)$. These translates have disjoint support interiors.
Indeed, if two sign patterns first differ at index $i$, their centers differ
by at least
\[
 2\left(a_i-\sum_{i<k\le r}a_k\right)
 \ge2\sum_{k>r}a_k,
\]
which is the diameter of a tail support. Their boundaries cause no additional
values because all tail densities are smooth and vanish at their support
endpoints.

The density of the tail starting at $r+1$ has supremum exactly
$1/(2a_{r+1})$. Its first uniform factor gives this upper bound, and at zero
it attains the bound because the rest of the tail is supported in
$[-a_{r+1},a_{r+1}]$. Taking the supremum of the disjoint translates proves
\eqref{eq:derivative-norms}. Finally
$\sum_{k>j}a_k\le a_j\sum_{k\ge1}2^{-k}=a_j$ under the ratio bound.
\end{proof}

For the dyadic reference this gives
\begin{equation}
 \norm{f_0^{(r)}}_\infty=2^{r(r+1)/2}.
 \label{eq:reference-norm}
\end{equation}
This special-case identity is also recorded in the repository's formal
crosswalk \cite{anchored}. We have supplied an independent proof of the
more general separated-support formula rather than assuming that crosswalk
as a formal verification of this manuscript.

\begin{definition}[The predecessor's smooth class at the boundary]
\label{def:smooth}
Let $\K$ consist of decreasing nonnegative spectra with
$\sum a_j\le3$, variance $1/9$, and a density $f_a\in C_c^\infty(\R)$
satisfying
\begin{equation}
 \norm{f_a^{(r)}}_\infty
 \le 2^{r+1}\norm{f_0^{(r)}}_\infty\qquad(r\ge0).
 \label{eq:smooth-class}
\end{equation}
This is the class used in the predecessor \cite{anchored}. Define
$\K_{1/2}=\K\cap\C_*$. The new notation extends its ratio-restricted
classes to the previously unresolved endpoint.
\end{definition}

\subsection{Delete one factor and restore the variance}

For $n\ge3$, put
\begin{equation}
 h_n=2^{-n},\qquad c_n=(1-3h_n^2)^{-1/2},\qquad
 a_j^{(n)}=
 \begin{cases}
 c_n2^{-j},&j<n,\\
 c_n2^{-(j+1)},&j\ge n.
 \end{cases}
 \label{eq:deleted-witness}
\end{equation}
This is the spectrum obtained by deleting the $n$th dyadic uniform factor,
reindexing, and rescaling all remaining factors to restore the variance.
Unlike a truncation, it keeps infinitely many positive factors and therefore
keeps a smooth density.

\begin{proposition}[Exact properties of the deleted-factor family]
\label{prop:deleted}
For every $n\ge3$, $a^{(n)}\in\K_{1/2}$. Writing $h=h_n$ and $c=c_n$,
\begin{align}
 \sum_j a_j^{(n)}&=c(1-h)\le1, \label{eq:witness-support}\\
 d_\infty(a^{(n)},a^0)
 &=\max\{(c-1)/2,\ h(1-c/2)\}\ge h/3, \label{eq:witness-gap}\\
 e_4(a^{(n)})
 &=\frac4{75}\frac{h^2(1-4h^2)}{(1-3h^2)^2}, \label{eq:witness-moment}\\
 \TV(\mu_{a^{(n)}},\mu_0)&\le4h^2. \label{eq:witness-tv}
\end{align}
In particular $d_\infty(a^{(n)},a^0)/h\to1/2$ and
$e_4(a^{(n)})/h^2\to4/75$.
\end{proposition}
\begin{proof}
All ratios are $1/2$ except at $n-1$, where the ratio is $1/4$.
The squared sum is $c^2(1/3-h^2)=1/3$. The support estimate follows from
$(1-h)^2\le1-3h^2$, valid for $h\le1/2$.

Before index $n$, the largest coefficient change is $(c-1)/2$.
From index $n$ onward, it is $h(1-c/2)$. Since $n\ge3$ gives
$c\le4/3$, this proves \eqref{eq:witness-gap}. The fourth-power sum is
$c^4(1/15-h^4)$, and direct subtraction yields
\[
 \sum_j(a_j^{(n)})^4-\frac1{15}
 =\frac25\frac{h^2(1-4h^2)}{(1-3h^2)^2}.
\]
Equation \eqref{eq:moments} proves \eqref{eq:witness-moment}.

For every derivative order $r$, the ratio of the first $r+1$ coefficient
products to the corresponding dyadic product is
$c^{r+1}2^{-\max(0,r-n+2)}$. By \cref{lem:derivative-norms},
\begin{equation}
 \frac{\norm{f_{a^{(n)}}^{(r)}}_\infty}{\norm{f_0^{(r)}}_\infty}
 =c^{-(r+1)}2^{\max(0,r-n+2)}\le2^{r+1}.
 \label{eq:witness-derivatives}
\end{equation}
This checks all the infinitely many inequalities in \eqref{eq:smooth-class},
not merely a finite set of derivative orders.

For the total variation estimate, let $v_n$ be the density of the dyadic
sum with the $n$th factor deleted. The original density is
$f_0=v_n*u_h$, where $u_h$ is the uniform density on $[-h,h]$.
The first two dyadic factors are still present. Distributional differentiation
of those factors, followed by convolution, gives
\begin{equation}
 \norm{v_n'}_1\le2,\qquad \norm{v_n''}_1\le8.
 \label{eq:common-derivatives}
\end{equation}
The centered Taylor remainder gives
\[
 \norm{v_n*u_h-v_n}_1\le\frac{h^2}{6}\norm{v_n''}_1,
 \qquad \TV(v_n*u_h,v_n)\le\frac23h^2.
\]
For the dilation operator $D_cv(x)=c^{-1}v(x/c)$, differentiating with
respect to $t=\log c$ and changing variables yields
\[
 \norm{D_cv-v}_1\le|\log c|\norm{v+xv'}_1.
\]
Since $v_n$ is supported in $[-1,1]$, \eqref{eq:common-derivatives} gives
$\norm{v_n+xv_n'}_1\le3$. Thus the additional total variation from
rescaling is at most $3\log c/2$. Finally
\[
 \frac23h^2+\frac32\log c
 \le h^2\left(\frac23+\frac{9}{4(1-3h^2)}\right)<4h^2
 \qquad(h\le1/8),
\]
using $-\log(1-x)\le x/(1-x)$. This proves \eqref{eq:witness-tv}.
The two limits follow from the explicit formulas.
\end{proof}

\begin{definition}[Anchored modulus]
For a class $\mathcal A$ containing $a^0$ and a law distance $D$, define
\begin{equation}
 \omega_{\mathcal A,D}(\eps)=
 \sup\{d_\infty(a,a^0):a\in\mathcal A,
                           \ D(\mu_a,\mu_0)\le\eps\}.
 \label{eq:modulus}
\end{equation}
\end{definition}

\begin{theorem}[Sharp boundary modulus in the exact smooth class]
\label{thm:sharp-modulus}
For $\mathcal A=\C_*$ or $\mathcal A=\K_{1/2}$, and all sufficiently
small $\eps>0$,
\begin{align}
 \frac1{12}\sqrt\eps
 &\le\omega_{\mathcal A,\TV}(\eps)
 \le\sqrt{\frac{75}{4}}\sqrt\eps, \label{eq:sharp-TV}\\
 \frac1{12}\sqrt\eps
 &\le\omega_{\mathcal A,\dK}(\eps)
 \le\sqrt{\frac{75}{2}}\sqrt\eps. \label{eq:sharp-K}
\end{align}
No exponent strictly greater than $1/2$ is valid in either class. The same
sharp exponent holds for every $\ell^p$ coefficient loss, $1\le p\le\infty$.
\end{theorem}
\begin{proof}
The upper estimates were proved in \cref{sec:recovery}. For the lower bound,
choose the smallest $n\ge3$ with $4h_n^2\le\eps$. For sufficiently small
$\eps$, minimality gives $16h_n^2>\eps$, hence $h_n>\sqrt\eps/4$.
By \cref{prop:deleted}, the law distance is at most $\eps$ and the spectral
gap is greater than $\sqrt\eps/12$. Kolmogorov distance is no larger than
total variation, so the same witnesses work there. The $\ell^p$ upper
bounds follow from \cref{prop:l1} and \eqref{eq:moment-metrics}; their lower
bounds follow from $\norm{a-a^0}_p\ge d_\infty(a,a^0)$.
\end{proof}

The constants in \eqref{eq:sharp-TV}--\eqref{eq:sharp-K} are not asserted
to be optimal. Sharpness here means the exact exponent and a matching rate
for every sufficiently small error level, not merely along a subsequence.

\section{A Hellinger estimate that sees the support boundary}
\label{sec:hellinger}

The total variation bound $O(h^2)$ alone gives only an $\Omega(h^{-2})$
testing lower bound by the elementary product-TV inequality. To obtain the
correct $h^{-4}$ scale, we need squared Hellinger distance of order $h^4$.
Because the laws have compact supports, a careless division by a density
outside its support would be invalid. The following lemmas explicitly handle
that boundary.

We use the convention
\begin{equation}
 H^2(f,g)=\int_\R(\sqrt f-\sqrt g)^2\dd x.
 \label{eq:hellinger-definition}
\end{equation}
Then $0\le H^2\le2$ and $\TV(f,g)\le H(f,g)$. Squared Hellinger distance
is jointly convex under common mixtures: Cauchy--Schwarz gives
\[
 \sqrt{\Bigl(\int f_z\dd\nu(z)\Bigr)
             \Bigl(\int g_z\dd\nu(z)\Bigr)}
 \ge\int\sqrt{f_zg_z}\dd\nu(z).
\]
Integration proves the claim, and convolution contraction is its special case
for common mixtures of translations.

\begin{lemma}[Centered smoothing costs fourth order in squared Hellinger]
\label{lem:hellinger-smoothing}
Let $v$ be the density of a sum of $m\ge5$ independent uniforms with fixed
positive half-lengths. There exist $C,h_0>0$ such that, for every centered
random variable $Z_h$ supported in $[-h,h]$, with $0<h<h_0$,
\begin{equation}
 H^2(v,v*\Law(Z_h))\le Ch^4.
 \label{eq:hellinger-smoothing}
\end{equation}
The constant is independent of the distribution of $Z_h$.
\end{lemma}
\begin{proof}
Write $B$ for the sum of the half-lengths, and $v_h=v*\Law(Z_h)$.
The convolution density $v$ is a $C^{m-2}$ piecewise polynomial, positive
on $(-B,B)$. For sufficiently small $s>0$ it has the exact endpoint form
\begin{equation}
 v(-B+s)=v(B-s)=C_0s^{m-1},
 \qquad
 C_0=\frac{1}{(m-1)!\prod_{i=1}^m(2\alpha_i)},
 \label{eq:spline-endpoint}
\end{equation}
where $\alpha_i$ are the half-lengths. This follows by convolving constant
densities on $[0,2\alpha_i]$: before reaching the smallest upper endpoint,
the convolution is the elementary simplex-volume polynomial.

On the interior $I_h=[-B+2h,B-2h]$, Taylor's theorem and $\E Z_h=0$ give
\begin{equation}
 |v_h(x)-v(x)|\le\frac{h^2}{2}
               \sup_{|u-x|\le h}|v''(u)|.
 \label{eq:smoothing-taylor}
\end{equation}
Also $(\sqrt{v_h}-\sqrt v)^2\le(v_h-v)^2/v$ there.
If $s$ is the distance to a nearby endpoint and $s\ge2h$, the endpoint
formula gives
\[
 \frac{\sup_{|u-x|\le h}|v''(u)|^2}{v(x)}
 \le C_1s^{m-5}.
\]
This is uniformly integrable because $m\ge5$. On a fixed central compact
interval, positivity of $v$ and boundedness of $v''$ give the same uniform
integral bound. Thus the Hellinger integral over $I_h$ is at most $C_2h^4$.

Outside $I_h$, both densities can be nonzero only within distance $3h$
of the old endpoints when their defining convolution arguments are considered.
Equation \eqref{eq:spline-endpoint} bounds the total mass of either density
in these boundary and exterior strips by $C_3h^m$. Since
$(\sqrt{v_h}-\sqrt v)^2\le v_h+v$, their contribution is
$O(h^m)=O(h^4)$. This also covers the part where $v=0$ and $v_h$ need
not vanish. Combining the two regions proves the result.
\end{proof}

\begin{lemma}[Uniform Hellinger control of small dilations]
\label{lem:hellinger-dilation}
For the density $v$ in \cref{lem:hellinger-smoothing} and fixed $R<\infty$,
there is $A_R<\infty$ such that every probability measure $\nu$ supported
in $[-R,R]$ and every $c>0$ satisfy
\begin{equation}
 H(D_c(v*\nu),v*\nu)\le A_R|\log c|.
 \label{eq:hellinger-dilation}
\end{equation}
\end{lemma}
\begin{proof}
First fix a translate $z\in[-R,R]$, and put $c=e^t$.
After translation by $-z$, comparison of $D_c(v(\cdot-z))$ with
$v(\cdot-z)$ becomes comparison of $v$ with the density of
$e^tY+(e^t-1)z$, where $Y$ has density $v$. Its square root is
\[
 s_t(x)=e^{-t/2}\sqrt{v(e^{-t}(x+z)-z)}.
\]
Differentiation and the substitution $u=e^{-t}(x+z)-z$ give
\begin{equation}
 \norm{\partial_ts_t}_2
 =\frac12\left(\int_{-B}^B
       \frac{(v(u)+(u+z)v'(u))^2}{v(u)}\dd u\right)^{1/2}.
 \label{eq:dilation-score}
\end{equation}
The integral is finite uniformly in $|z|\le R$: near an endpoint the
potentially singular term $v'(u)^2/v(u)$ is $O(s^{m-3})$, and on central
compact intervals $v$ is bounded below. There are no distributional boundary
terms for the square root. Indeed \eqref{eq:spline-endpoint} shows that
$\sqrt v$, extended by zero, lies in $H^1(\R)$ and vanishes at the endpoints.
The displayed differentiation is justified first for smooth compactly supported
approximations to $\sqrt v$ and then in $L^2$, using convergence of the
function and its derivative; the supports are uniformly bounded.

The right side of \eqref{eq:dilation-score} is independent of $t$ and
bounded by some $A_R$. The fundamental theorem of calculus in $L^2$
therefore gives $\norm{s_t-s_0}_2\le A_R|t|$.
Finally mix over the same $z\sim\nu$ and use joint convexity of $H^2$.
\end{proof}

\begin{theorem}[Hellinger order of the deleted-factor witnesses]
\label{thm:witness-hellinger}
There are constants $0<c_H\le C_H<\infty$ such that, for every $n\ge6$,
\begin{equation}
 c_Hh_n^4\le H^2(\mu_{a^{(n)}},\mu_0)\le C_Hh_n^4.
 \label{eq:witness-hellinger}
\end{equation}
\end{theorem}
\begin{proof}
Let $v$ be the density of the first five dyadic uniforms. The density $v_n$
of the sum with the $n$th factor deleted has the form $v*\nu_n$, with
$\nu_n$ supported in a fixed bounded interval. By convolution contraction
and \cref{lem:hellinger-smoothing},
\[
 H(v_n,v_n*u_{h_n})\le H(v,v*u_{h_n})\le C h_n^2.
\]
By \cref{lem:hellinger-dilation},
$H(D_{c_n}v_n,v_n)\le A\log c_n=O(h_n^2)$.
Since $v_n*u_{h_n}=f_0$, the triangle inequality proves the upper bound
for all sufficiently large $n$. Increasing the constant covers the finitely
many remaining indices $n\ge6$.

For the lower bound, \eqref{eq:moment-metrics} and $\TV\le H$ imply
$H^2\ge e_4(a^{(n)})^2$. Formula \eqref{eq:witness-moment} has a positive
ratio to $h_n^2$ tending to $4/75$, so it is bounded below by a positive
constant times $h_n^2$ on all $n\ge6$. Squaring finishes the proof.
\end{proof}

The proof only needs five unchanged positive uniform factors, centered small
noise, and a bounded remaining support. This is why it also extends to
nongeometric spectra in \cref{sec:profiles}. It does not require positivity
of a compactly supported density on the entire real line, finite chi-square
distance across unequal supports, or an unproved pointwise likelihood expansion.

\section{Quartic sample complexity and honest confidence sets}
\label{sec:statistics}

\subsection{A self-contained concentration input}

For independent $Z_i\in[0,1]$ with common mean $m$, the standard Hoeffding
bound \cite{hoeffding} is
\begin{equation}
 \PP\{\overline Z_N-m\ge t\}\le e^{-2Nt^2},\qquad
 \PP\{m-\overline Z_N\ge t\}\le e^{-2Nt^2}.
 \label{eq:hoeffding}
\end{equation}
For completeness, if
$\psi(\lambda)=\log\E e^{\lambda(Z-\E Z)}$, then $\psi(0)=\psi'(0)=0$
and $\psi''(\lambda)$ is a variance under an exponentially tilted law on
$[0,1]$. That variance is at most $1/4$, because
$\Var Z\le\E(Z-1/2)^2\le1/4$. Hence
$\psi(\lambda)\le\lambda^2/8$. Independence, Markov's inequality,
and optimization at $\lambda=4t$ prove the first bound for the average;
apply it to $1-Z_i$ for the second. This is a proof of the standard input,
not a new concentration inequality.

\subsection{The testing theorem}

For $\mathcal A=\C_*$ or $\K_{1/2}$, let $N_{\mathcal A}(\delta)$ be the
smallest positive integer $N$ admitting a possibly randomized test $\psi$
with
\begin{equation}
 \E_{\mu_0^{\otimes N}}\psi\le\tfrac14,
 \qquad
 \sup_{\substack{a\in\mathcal A\\d_\infty(a,a^0)\ge\delta}}
 \E_{\mu_a^{\otimes N}}(1-\psi)\le\tfrac14.
 \label{eq:testing-definition}
\end{equation}

\begin{theorem}[Sharp quartic testing complexity]
\label{thm:testing}
For either class, $N_{\mathcal A}(\delta)=\Theta(\delta^{-4})$ as
$\delta\downarrow0$. A concrete upper bound is
\begin{equation}
 N_{\mathcal A}(\delta)\le
 \left\lceil\frac{5625\log4}{8\delta^4}\right\rceil.
 \label{eq:sample-upper}
\end{equation}
\end{theorem}
\begin{proof}
Given observations $X_1,\ldots,X_N$, set
$\overline Z_N=N^{-1}\sum_{i=1}^NX_i^4$ and reject the null if
\begin{equation}
 \frac{19}{675}-\overline Z_N>\frac{2\delta^2}{75}.
 \label{eq:fourth-test}
\end{equation}
Every $X_i$ lies in $[-1,1]$, so $Z_i=X_i^4\in[0,1]$.
Under the null, \eqref{eq:hoeffding} bounds rejection probability by
$\exp(-8N\delta^4/5625)$. Any alternative in
\eqref{eq:testing-definition} has $e_4(a)\ge4\delta^2/75$ by
\eqref{eq:dyadic-moment-main}. Acceptance then requires an upward empirical
moment error at least $2\delta^2/75$, with the same probability bound.
This proves \eqref{eq:sample-upper}.

For the lower bound, choose $n$ so that $3\delta\le h_n<6\delta$;
for small $\delta$ we have $n\ge6$. By \cref{prop:deleted}, this is an
alternative in $\K_{1/2}$ at spectral distance at least $\delta$.
For any probability laws $P,Q$, multiplicativity of Hellinger affinity gives
\begin{equation}
 H^2(P^{\otimes N},Q^{\otimes N})
 =2\left[1-\left(1-\tfrac12H^2(P,Q)\right)^N\right]
 \le NH^2(P,Q).
 \label{eq:product-hellinger}
\end{equation}
Combining \cref{thm:witness-hellinger}, \eqref{eq:product-hellinger}, and
$\TV\le H$ yields
\[
 \TV(\mu_{a^{(n)}}^{\otimes N},\mu_0^{\otimes N})
 \le\sqrt{NC_Hh_n^4}\le\sqrt{NC_H6^4\delta^4}.
\]
For $N\le(16C_H6^4)^{-1}\delta^{-4}$ this is at most $1/4$.
The two errors of any test at these two laws sum to at least one minus
this total variation, hence at least $3/4$. They cannot both be at most
$1/4$. This proves the lower order for $\K_{1/2}$ and therefore for
$\C_*$ as well.
\end{proof}

The constant in the lower bound is not numerically optimized. Its finiteness
is established analytically by the explicit spline endpoint and dilation-score
integrals. No numerical Hellinger integration is being used as a proof.

\subsection{Confidence sets and an observable radius}

For $0<\alpha<1$, set
\begin{equation}
 t_N=\sqrt{\frac{\log(2/\alpha)}{2N}},\qquad
 \mathcal F_N=\{a\in\C_*:|m_4(a)-\overline Z_N|\le t_N\}.
 \label{eq:confidence}
\end{equation}
By the two tails of \eqref{eq:hoeffding},
$\PP_a\{a\in\mathcal F_N\}\ge1-\alpha$ uniformly over $\C_*$.
At the reference, on an event of probability at least $1-\alpha$, every
candidate in $\mathcal F_N$ satisfies $e_4(a)\le2t_N$. Hence
\begin{align}
 \sup_{a\in\mathcal F_N}d_\infty(a,a^0)
 &\le\sqrt{\frac{75}{2}t_N}, \label{eq:confidence-global}\\
 \sup_{a\in\mathcal F_N}d_r(a,a^0)
 &\le\frac{75}{2}2^r t_N. \label{eq:confidence-prefix}
\end{align}
On that event $a^0\in\mathcal F_N$, so no convention about the supremum
of an empty confidence set is needed. This establishes honest coverage and
reference-point contraction, not a finite representation of the entire
infinite-dimensional set. The scalar radius and test themselves are directly
computable from a fourth-moment average.

\section{Sharp local conditioning of a growing prefix}
\label{sec:prefix}

For $D=\TV$ or $\dK$ and $\mathcal A=\C_*$ or $\K_{1/2}$, define
\begin{equation}
 \Lambda_r^D(\mathcal A)=\lim_{\eps\downarrow0}
 \sup_{\substack{a\in\mathcal A\\0<D(\mu_a,\mu_0)\le\eps}}
 \frac{d_r(a,a^0)}{D(\mu_a,\mu_0)}.
 \label{eq:prefix-constant}
\end{equation}
The limit exists as a decreasing limit of suprema, allowing infinity in
principle. The earlier upper bounds make it finite here.

\begin{theorem}[The prefix constant grows exactly geometrically]
\label{thm:prefix}
For every $r\ge6$ and either class,
\begin{align}
 \frac{2^r}{30}&\le\Lambda_r^{\TV}(\mathcal A)
                       \le\frac{75}{4}2^r, \label{eq:prefix-TV}\\
 \frac{2^r}{30}&\le\Lambda_r^{\dK}(\mathcal A)
                       \le\frac{75}{2}2^r. \label{eq:prefix-K}
\end{align}
\end{theorem}
\begin{proof}
The upper estimates follow directly from \eqref{eq:dyadic-prefix-main}
and \eqref{eq:moment-metrics}. For the lower estimates set $n=r$, $h=2^{-n}$,
and, for $1/2\le\theta<1$, put
\begin{equation}
 \lambda=1-\theta^2,\qquad c=(1-4\lambda h^2)^{-1/2},\qquad
 a_j(\theta)=
 \begin{cases}
 c2^{-j},&j<n,\\
 c\theta2^{-j},&j\ge n.
 \end{cases}
 \label{eq:tail-deformation}
\end{equation}
The dyadic tail starting at $n$ has squared coefficient sum $4h^2/3$.
Thus the chosen $c$ restores squared sum $1/3$. The only altered ratio is
at $n-1$, where it is $\theta/2\le1/2$.
For every derivative order $k$, \cref{lem:derivative-norms} gives the norm
ratio
\[
 c^{-(k+1)}\theta^{-\max(0,k-n+2)}\le2^{k+1}.
\]
Together with \cref{thm:support}, this proves $a(\theta)\in\K_{1/2}$.

At coordinate $n$, $c\theta<1$, and
\begin{equation}
 1-c\theta
 =\frac{\lambda(1-4h^2)}{(1-4\lambda h^2)(1+c\theta)}
 \ge\frac{1-\theta}{3}.
 \label{eq:theta-gap}
\end{equation}
Here the denominator is at most $2$, $\lambda\ge1-\theta$, and
$h\le1/8$. Consequently $d_n(a(\theta),a^0)\ge(1-\theta)h/3$.

We next show
\begin{equation}
 \TV(\mu_{a(\theta)},\mu_0)\le10(1-\theta)h^2.
 \label{eq:theta-TV}
\end{equation}
Let $w$ be the density of the head before $n$, and let $T$ be the independent
dyadic random tail starting at $n$. Then $\E T=0$ and
$\E T^2=4h^2/9$. For $s\in[\theta,1]$, differentiating
$\E w(x-sT)$ and subtracting $\E[T w'(x)]=0$ yields
\[
 \norm{\tfrac{\mathrm d}{\mathrm ds}\E w(\cdot-sT)}_1
 \le s\E T^2\norm{w''}_1.
\]
Integration gives an $L^1$ difference at most
$\lambda\E T^2\norm{w''}_1/2$. Since the head contains the first five
factors, $\norm{w''}_1\le8$. The resulting TV bound is
$8\lambda h^2/9$. The law before rescaling retains a first uniform factor
of half-length $1/2$, has derivative $L^1$ norm at most $2$, and is supported
in $[-1,1]$. Its rescaling costs TV at most
\[
 \frac32\log c\le\frac{3\lambda h^2}{1-4\lambda h^2}
 \le4\lambda h^2.
\]
The total is at most $5\lambda h^2\le10(1-\theta)h^2$, proving
\eqref{eq:theta-TV}.

For $\theta<1$ the spectrum is not geometric, so moment rigidity ensures
that both law distances are positive. As $\theta\uparrow1$, they tend to
zero by \eqref{eq:theta-TV}. Dividing \eqref{eq:theta-gap} times $h$ by
\eqref{eq:theta-TV} gives a limiting ratio at least $1/(30h)=2^r/30$.
The same lower bound applies to $\dK\le\TV$. Since these witnesses are
in $\K_{1/2}$, it applies to both classes.
\end{proof}

This conditioning law is substantially more favorable than the
$\exp(O(r^2))$ upper constants in the predecessor's prefix theorem with an
unrestricted infinite tail \cite{anchored}. It is a comparison between
\emph{different constraint classes}, not a correction to that theorem.
The exact critical constraint supplies a new order structure that the larger
class does not possess.

\section{Nongeometric reference spectra and relative-ratio cones}
\label{sec:profiles}

The geometric zero lattice is not needed for the rigidity argument. The
appropriate general structure is monotonicity of the coefficient ratio to a
fixed reference, which translates into a monotone likelihood ratio between
two probability sequences.

Let $b=(b_j)_{j\ge1}$ be positive with
\begin{equation}
 \sum_jb_j^2=S>0,\qquad
 0<r_-\le\frac{b_{j+1}}{b_j}\le r_+<1.
 \label{eq:profile-ratios}
\end{equation}
Define the relative-ratio cone
\begin{equation}
 \C_b=\left\{a_j\ge0:\sum_ja_j^2=S,\quad
       \frac{a_{j+1}}{b_{j+1}}\le\frac{a_j}{b_j}\quad(j\ge1)\right\}.
 \label{eq:profile-class}
\end{equation}
In particular $a_{j+1}\le r_+a_j$, so there is a common support bound
$\sum a_j\le\sqrt S/(1-r_+)$. Write
\[
 w_j=b_j^2/S,\quad p_j=a_j^2/S,\quad
 \Delta_b=\sum_jp_j^2-\sum_jw_j^2.
\]

\begin{theorem}[Relative-ratio moment rigidity]
\label{thm:profile-rigidity}
For every $a\in\C_b$,
\begin{equation}
 \Delta_b\ge0,\qquad
 d_\infty(a,b)^2
 \le\frac{S\Delta_b}{2w_1(w_1-w_2)}
 =\frac{15\,[m_4(b)-m_4(a)]}{4S w_1(w_1-w_2)}.
 \label{eq:profile-bound}
\end{equation}
Equality of the second and fourth moments with those of $b$ forces $a=b$.
For this upper conclusion, the lower ratio bound $r_->0$ is unnecessary;
it suffices that the positive reference has $w_1>w_2$, is decreasing, and
is summable, with a common support bound for the metric consequence.
\end{theorem}
\begin{proof}
The sequence $p_j/w_j=(a_j/b_j)^2$ is nonincreasing, with $w$-weighted
average one. Its weighted average on any leading block is at least its
weighted average on the complementary tail, and therefore is at least one.
Thus $D_k=\sum_{j\le k}(p_j-w_j)\ge0$. Summation by parts, exactly as
in \cref{lem:energy}, gives
\begin{equation}
 \Delta_b=\sum_j(p_j-w_j)^2
         +2\sum_{k\ge1}(w_k-w_{k+1})D_k.
 \label{eq:profile-energy}
\end{equation}
The terms are nonnegative and the sums converge. Set $d=p_1-w_1=D_1$.
Then $\Delta_b\ge2(w_1-w_2)d$.

Since $p_j/w_j\le p_1/w_1=1+d/w_1$, the sum of the positive differences
$p_j-w_j$ is at most $d/w_1$. Its value equals the absolute sum of the
negative differences because both sequences have total mass one. Hence
$|p_j-w_j|\le d/w_1$ at every coordinate. Taking square roots as in
\cref{thm:geometric-rigidity} proves the first bound in
\eqref{eq:profile-bound}. The moment formula \eqref{eq:moments} proves
the second. If the moment deficit is zero, the first term of
\eqref{eq:profile-energy} vanishes and $p=w$, hence $a=b$.
\end{proof}

\begin{theorem}[Sharp recovery and testing for a fixed ratio profile]
\label{thm:profile-sharp}
Under \eqref{eq:profile-ratios}, the modulus anchored at $b$ over $\C_b$
is $\Theta(\sqrt\eps)$ for either TV or Kolmogorov distance. Testing
$b$ against alternatives $a\in\C_b$ with $d_\infty(a,b)\ge\delta$ has
sample complexity $\Theta(\delta^{-4})$ at fixed errors $1/4$.
The lower witnesses have infinitely many positive coefficients and compactly
supported $C^\infty$ densities.
\end{theorem}
\begin{proof}
The upper modulus follows from \cref{thm:profile-rigidity} and
\eqref{eq:moment-metrics} with the common support bound just stated.
For lower witnesses, fix $\theta=1/2$, and let
\begin{equation}
 \tau_n=\sum_{j\ge n}b_j^2,\qquad
 c_n=\left(1-(1-\theta^2)\tau_n/S\right)^{-1/2},\qquad
 a_j^{[n]}=
 \begin{cases}c_nb_j,&j<n,\\c_n\theta b_j,&j\ge n.\end{cases}
 \label{eq:profile-witness}
\end{equation}
These spectra belong to $\C_b$, and
\[
 b_n^2\le\tau_n\le\frac{b_n^2}{1-r_+^2},\qquad
 c_n-1=O_b(b_n^2),\qquad
 d_\infty(a^{[n]},b)\ge b_n/4
\]
for all sufficiently large $n$. The last inequality follows at coordinate
$n$ from $c_n\theta\to1/2$.

Use the first five reference uniforms as the fixed convolution density $v$
in \cref{lem:hellinger-smoothing}. The remaining head before $n$ supplies
a common convolution factor. The centered tail starting at $n$ is supported
in an interval of radius at most $b_n/(1-r_+)$, and its contracted version
has the same type of bound. Comparing each with the unchanged head by
\cref{lem:hellinger-smoothing} and then using the triangle inequality gives
Hellinger distance $O_b(b_n^2)$ between the unscaled laws. Rescaling adds
$O_b(\log c_n)=O_b(b_n^2)$ by \cref{lem:hellinger-dilation}, whose residual
support bound is uniform. Therefore
\begin{equation}
 H^2(\mu_{a^{[n]}},\mu_b)=O_b(b_n^4),\qquad
 \TV(\mu_{a^{[n]}},\mu_b)=O_b(b_n^2).
 \label{eq:profile-H}
\end{equation}
Positivity of every coefficient gives smoothness by the Fourier-decay
argument in \cref{lem:derivative-norms}; separated translates are not needed
for this part.

The lower ratio bound $b_{n+1}\ge r_-b_n$ lets us choose $b_n$ within a
fixed multiplicative factor of any sufficiently small target scale. Thus
\eqref{eq:profile-H} and the spectral gap imply a lower modulus of order
$\sqrt\eps$ for every sufficiently small $\eps$, not just selected errors.
For testing, the fourth-moment deficit in \eqref{eq:profile-bound} is at
least a fixed multiple of $\delta^2$, so the bounded-variable argument in
\cref{thm:testing} gives an $O_b(\delta^{-4})$ upper bound. Choose $n$
with $4\delta\le b_n<4\delta/r_-$. Equations
\eqref{eq:profile-H} and \eqref{eq:product-hellinger} give the matching
lower bound just as before.
\end{proof}

Every geometric reference in \eqref{eq:general-class} satisfies
\eqref{eq:profile-ratios} with $r_-=r_+=\rho$, and its relative-ratio
cone is exactly $\C_{\rho,S}$. Thus the sharp exponent and quartic testing
law hold for every geometric reference, not only the dyadic one. For those
geometric references, \cref{prop:l1} supplies the matching upper bounds for
all $\ell^p$ losses; the same witnesses supply the lower bounds.

Without a positive lower ratio bound, \cref{thm:profile-rigidity} still gives
its upper estimate, and \eqref{eq:profile-H} still obstructs every better
H\"older exponent along a sequence. What may fail is interpolation between
widely separated scales. No all-error matching modulus is claimed for that
more general case.

\section{What has been resolved, and what has not}
\label{sec:scope}

\begin{center}\small
\begin{tabular}{@{}>{\raggedright\arraybackslash}p{0.24\textwidth}
>{\raggedright\arraybackslash}p{0.66\textwidth}@{}}
\toprule
Question or regime & Conclusion and precise scope\\
\midrule
Exact dyadic ratio $1/2$ & Sharp square-root anchored modulus, including the
predecessor's exact smooth class; \cref{thm:sharp-modulus}.\\[5pt]
Fixed variance and two moments & The dyadic and geometric moment fibres at
the reference are singletons inside their critical cones;
\cref{thm:geometric-rigidity}.\\[5pt]
Exact maximal support & With the critical ratio and fixed variance, fixing
the reference support endpoint forces the reference itself;
\cref{thm:support}.\\[5pt]
Leading $r$ coefficients & Local conditioning is $\Theta(2^r)$ in the
critical dyadic cone, not uniformly over unrestricted tails;
\cref{thm:prefix}.\\[5pt]
Statistical complexity & Anchored testing is $\Theta(\delta^{-4})$; honest
reference contraction is $O(N^{-1/4})$;
\cref{thm:testing}.\\[5pt]
Looser dyadic ratio bound & The predecessor's fixed-slack stretched-exponential
result is cited, not reproved or newly formalized here.\\[5pt]
Arbitrary two-point inverse & Not settled. The moment extremizer is a fixed
anchor; the proofs do not compare two free spectra.\\
\bottomrule
\end{tabular}
\end{center}

The apparent abrupt transition at $1/2$ has a precise cause. At the exact
boundary, normalized squared coefficients are forced to majorize the
reference sequence. Their fourth-power deficit cannot be canceled. With
slack, that one-sided relation to the dyadic reference disappears, permitting
the moment-matching perturbations investigated by the predecessor. This
explanation is supported by \eqref{eq:energy}, rather than by numerical
conditioning experiments.

The results are not claimed as a resolution of a famous general inverse
problem. They resolve a specific, explicitly stated research frontier in the
repository and give additional quantitative consequences with complete proofs.
Their relevance is the contrast between severe infinite-dimensional
instability and a two-moment certificate under an exact structural constraint.

\section{Verification, provenance, and formalization route}
\label{sec:audit}

\subsection{Exact computations and their limits}

The accompanying standard-library Python program \path{code/verify.py}
uses exact rational arithmetic for finite regression tests and 65-digit
\texttt{Decimal} arithmetic only for illustrative radical evaluations.
The recorded run used Python 3.13.5, deterministic seed 20260930, and passed
\textbf{29,776 exact assertions}. The detailed counts are in
\path{results/verification.json}. They include the infinite-tail energy
identity evaluated symbolically for finite padded sequences, the sharp head
inequality, squared-coefficient bounds, exact fourth-moment formulas,
derivative-class inequalities for tested orders, tail-deformation estimates,
and nongeometric finite-profile identities.

These are finite regression checks. In particular, checking derivative orders
$0\le r\le100$ is \emph{not} evidence for all orders by itself; the universal
proof is \eqref{eq:witness-derivatives}. The Hellinger lemmas and infinite
limiting assertions are proved in the text, not certified by numerical tests.
The program performs no empirical integration of TV or Hellinger distance.

\begin{table}[htbp]
\centering\small
\begin{tabular}{@{}rrrr@{}}
\toprule
$n$ & Spectral gap & Fourth-moment deficit & TV upper bound $4h_n^2$\\
\midrule
3 & $6.0981560\cdot10^{-2}$ & $8.5998388\cdot10^{-4}$ & $6.2500000\cdot10^{-2}$\\
4 & $3.1065269\cdot10^{-2}$ & $2.0997047\cdot10^{-4}$ & $1.5625000\cdot10^{-2}$\\
6 & $7.8096374\cdot10^{-3}$ & $1.3027193\cdot10^{-5}$ & $9.7656250\cdot10^{-4}$\\
8 & $1.9530803\cdot10^{-3}$ & $8.1382692\cdot10^{-7}$ & $6.1035156\cdot10^{-5}$\\
12 & $1.2207030\cdot10^{-4}$ & $3.1789148\cdot10^{-9}$ & $2.3841858\cdot10^{-7}$\\
\bottomrule
\end{tabular}
\caption{Rounded evaluations of \cref{prop:deleted}. The last column is a
proved upper bound, not the actual TV distance. Decimal rounding is not
outward-directed interval certification.}
\label{tab:witnesses}
\end{table}

Run \texttt{python code/verify.py} from the package root, or by absolute
path from any directory. By default the outputs go to \path{results-rerun/},
so the recorded results are preserved. The option
\texttt{-{}-output-dir PATH} selects another destination. The TeX source
contains its tables directly and does not require Python or network access
to compile. Run \texttt{latexmk -pdf article.tex}, or run
\texttt{pdflatex article.tex} repeatedly until cross-references stabilize.
No shell escape is required.

\subsection{Pinned repository source and question mapping}

The source inspection used the repository snapshot
\begin{quote}\small\ttfamily\raggedright\repoSHA\end{quote}
and the article at
\begin{quote}\small\ttfamily\raggedright
Analysis/FabiusFunction/docs/\allowbreak
semi-formalized-research-frontiers/\allowbreak
drafts/inverse-and-sampling/\allowbreak
Anchored\_Dyadic\_Recovery\_Uniform\_Spectrum/\allowbreak article.tex.
\end{quote}
Its content blob was
\begin{quote}\small\ttfamily\raggedright\blobSHA.\end{quote}
The opening definitions, the smooth-class and derivative sections, the
statistical section, and the further-research questions were read through
the GitHub connector. The exact target is the question titled
\emph{The sharp separation boundary $\rho=1/2$}, third in the final question
list. The immediately following exact-support question is addressed here
only at the critical boundary. The earlier article's general identifiability,
Fourier zero-count arguments, and numerical bounds are not assumptions in
our new critical-boundary proofs.
% ed. (2026-09-30): reciprocal note, and related repository results the article does not cite
\begin{ednote}
The question answered here,
in \path{../Anchored_Dyadic_Recovery_Uniform_Spectrum/}, now carries a
note pointing to this package. At the dyadic reference this also settles
the boundary case $q=1/2$ that the editorial note to ``Geometrically
separated classes'' in \path{../Recovering_Uniform_Factors_Fabius_Rvachev/}
leaves open. In the same group,
\path{../Sharp_Stability_Strata_Fabius_Rvachev_Deconvolution/} proves, for
a finite list of uniform factors, that the optimal anchored exponent is
$1/2$ whenever the reference has a repeated positive scale or a zero slot
(its \texttt{thm:anchored}); the $\Theta(\sqrt\eps)$ modulus of
\cref{thm:sharp-modulus} is an infinite-spectrum counterpart at the dyadic
reference, proved independently.
\end{ednote}

The source's editorial crosswalk reports the formal derivative-norm theorem
\path{Fabius.isGreatest_abs_iteratedDeriv_rvachevUp} in
\path{GlobalBounds.lean} and the variance theorem
\path{Fabius.integral_sq_mul_rvachev_eq_one_ninth} in
\path{DyadicSpecializations.lean}, under
\path{Analysis/FabiusFunction/Lean/FabiusFunction/}.
These names are recorded as repository provenance. No Lean build was run for
this article, and no claim is made that the newly stated stability,
Hellinger, or statistical theorems are already formalized.

The literature check was targeted: Arias de Reyna for the compact-support
uniform-sum function, Billey--Swanson for generalized uniform-sum
identifiability, Hoeffding for concentration, and repository search for
related inverse and lacunarity reports. This is not an exhaustive novelty
audit. No repository files were modified.

\subsection{Claim-level dependency audit}

\begin{center}\small
\begin{tabular}{@{}>{\raggedright\arraybackslash}p{0.28\textwidth}
>{\raggedright\arraybackslash}p{0.62\textwidth}@{}}
\toprule
Conclusion & Mathematical dependencies\\
\midrule
Two-moment certificate & Tail domination, summation by parts, the sharp
head inequality, envelope slacks, and the elementary fourth-moment formula.\\[4pt]
Exact support rigidity & The same majorization inequality applied to
normalized coefficients rather than their squares.\\[4pt]
Smooth lower witnesses & Exact variance rescaling, separated translates of
uniform derivatives, and an $L^1$ centered Taylor remainder.\\[4pt]
Fourth-order Hellinger & Fixed finite uniform convolution, endpoint simplex
polynomials, separate boundary integration, and a square-root dilation score.\\[4pt]
Sharp testing order & Fourth-moment exposure, bounded-variable concentration,
the explicit witnesses, and product Hellinger affinity.\\[4pt]
Sharp prefix constants & Global coordinate certificate and a continuous
variance-preserving contraction of an entire tail.\\[4pt]
Nongeometric extension & Monotone relative ratios, a general weighted
summation-by-parts identity, and the same smoothing mechanism.\\
\bottomrule
\end{tabular}
\end{center}

A practical formalization route begins with finite partial-sum identities and
nonnegative geometric tails. One can then formalize the probability-sequence
energy identity, the quadratic head estimate, and the envelope-slack argument
without any Fabius-specific analysis. The fourth-moment identity is a separate
finite-independence lemma followed by dominated convergence. These modules
already give the main upper certificate and rigidity theorem.

The analytic lower-bound modules should remain separate: convolution
smoothness; exact derivative norms from disjoint support interiors; $L^1$
smoothing and dilation; then the endpoint-sensitive Hellinger lemmas.
The latter require an explicit zero-extension and Sobolev differentiation
interface, rather than a blanket differentiability assertion for likelihood
ratios. Finally, the testing theorem is assembled from deterministic moment
exposure and product-measure inequalities. This is a dependency plan, not a
claim about the availability of every needed theorem in a current mathlib API.

\section{Further research questions}
\label{sec:further}

\begin{question}[Optimal constants and possible log-periodic structure]
Do the limits $\omega_{\C_*,D}(\eps)/\sqrt\eps$ exist for $D=\TV$ and
$\dK$? The present constants are deliberately nonoptimal. The deleted-factor
scales are geometric, while continuous tail contractions provide an extra
parameter. Determine whether optimal combinations eliminate log-periodic
oscillation, and whether the explicit moment constant $75/4$ in the certificate
can itself be improved. Sharpness of the exponent does not answer either
constant question.
\end{question}

\begin{question}[A joint small-slack crossover]
Let the allowed dyadic ratio be $1/2+\eta$ with $\eta\downarrow0$ as the
law error $\eps\downarrow0$. Determine the two-parameter inverse modulus.
The critical square-root theorem and the predecessor's fixed-positive-slack
stretched-exponential theorem concern different orders of limits. A
quantitative measure of the negative parts of the cumulative discrepancies
in \eqref{eq:energy} may interpolate between the regimes.
\end{question}

\begin{question}[Pairwise recovery within the critical cone]
What is the optimal modulus between two freely varying spectra in $\C_*$
near $a^0$? The fourth moment exposes the reference, but its deficits for two
unknown spectra can be nearly equal. The anchored certificate therefore does
not automatically give a pairwise inverse theorem. One should distinguish
an estimator's uniform risk on a neighborhood from confidence-set contraction
at its extremal point.
\end{question}

\begin{question}[Ratio profiles without uniform scale spacing]
Classify all-error inverse moduli for relative-ratio cones when
$b_{j+1}/b_j\to0$ or $b_{j+1}/b_j\to1$. The algebraic upper certificate
survives far beyond \eqref{eq:profile-ratios}, but its rate need not be
attained at every error level by the sparse tail scales. A profile-dependent
modulus may retain more information than a single H\"older exponent.
\end{question}

\begin{question}[Exact support and variance with supercritical slack]
At the exact boundary the maximal-support fibre is a singleton. With a fixed
ratio allowance $\rho>1/2$, can the predecessor's very small law perturbations
be made to preserve both variance and the exact support endpoint? The present
rigidity theorem identifies why the critical version is impossible, but does
not settle the construction under slack.
\end{question}

\begin{question}[Other factor laws and Gaussian nuisance components]
For which centered base laws does a fourth-cumulant extremizer produce the
same square-root inverse law? The sign and nonvanishing of the fourth cumulant
matter for moment exposure. An additional unknown Gaussian variance changes
the second-moment normalization while leaving the fourth cumulant unchanged;
separate identifiability and stability statements are needed rather than an
unqualified transfer of the present theorem.
\end{question}

\begin{question}[Sharper statistical constants and robust measurements]
Determine local testing constants through the leading square-root-density
perturbation, and compare the fourth-moment test with a likelihood-based test.
Extend the confidence certificate to validated moment intervals, bounded
measurement error, or controlled contamination. Such extensions must keep
track of the exact structural constraint; an unverified ratio assumption is
not supplied by a good empirical fit to four moments.
\end{question}

\begin{question}[A formalized moment-exposure library]
Formalize the weighted identity \eqref{eq:profile-energy} and its quantitative
consequences as a reusable library for monotone relative-ratio cones.
A useful milestone is a kernel-checked proof of the dyadic constant in
\eqref{eq:dyadic-moment-main}, followed by the support singleton theorem.
The Hellinger boundary module is a separate, more analytic milestone.
\end{question}

\section{Conclusion}

The exact dyadic separation boundary is not another instance of the
fixed-slack instability. It has a rigid moment geometry. A fourth-moment
deficit controls the squared error of every coefficient, producing a sharp
square-root inverse law for the entire infinite spectrum and a sharp
geometric conditioning law for finite prefixes. Smooth deleted-factor
witnesses show that these estimates cannot be improved in exponent, even
under the predecessor's explicit derivative bounds. A separate Hellinger
argument yields the matching quartic testing complexity.

The essential mechanism is not peculiar to a dyadic Fourier zero lattice.
It is the one-sided majorization induced by relative-ratio constraints,
which also gives rigidity and sharp rates around nongeometric references.
The unresolved questions now lie in optimal constants, two-point rather
than anchored recovery, and the interaction of small structural slack with
small observational error.

\clearpage
\begin{thebibliography}{9}
\raggedright

\bibitem{anchored}
ProveIt research collection,
\emph{Anchored Recovery of the Dyadic Uniform Spectrum: A sharp
stretched-exponential modulus, lacunary counterexamples to H\"older stability,
and Lipschitz recovery of every fixed prefix},
research manuscript prepared for Vladimir Reshetnikov, 29 September 2026.
Inspected at commit \texttt{\repoSHA}; exact path and blob in
\cref{sec:audit}. Repository research article, not a refereed publication.
\url{https://github.com/VladimirReshetnikov/ProveIt/blob/f4457a1d19701d32e8f324240c497653d235b53e/Analysis/FabiusFunction/docs/semi-formalized-research-frontiers/drafts/inverse-and-sampling/Anchored_Dyadic_Recovery_Uniform_Spectrum/article.tex}.

\bibitem{arias}
J.~Arias de Reyna,
\emph{An infinitely differentiable function with compact support:
Definition and properties}, arXiv:1702.05442, 2017.
English translation of \emph{Definici\'on y estudio de una funci\'on
indefinidamente diferenciable de soporte compacto},
Rev. Real Acad. Ciencias Madrid \textbf{76} (1982), 21--38.
\url{https://arxiv.org/abs/1702.05442}.

\bibitem{billey}
S.~C.~Billey and J.~P.~Swanson,
\emph{The metric space of limit laws for $q$-hook formulas},
2022. DOI: \href{https://doi.org/10.5070/C62257868}{10.5070/C62257868}.
Author manuscript arXiv:2010.12701, version 2 (2023).
\url{https://arxiv.org/abs/2010.12701}.

\bibitem{hoeffding}
W.~Hoeffding,
\emph{Probability inequalities for sums of bounded random variables},
Journal of the American Statistical Association \textbf{58} (301)
(1963), 13--30.
DOI: \href{https://doi.org/10.1080/01621459.1963.10500830}
{10.1080/01621459.1963.10500830}.

\end{thebibliography}
\end{document}
